This work investigated how effectively synthetic images can be used to improve the performance of an object detection model. The object detection task of the work was to detect people in small indoor environment using images captured with a near-infrared sensor. Based on a literature review, different methods for generating synthetic images were compared, including 3D simulators, image blending, image-to-image and text-to-image models, for enhancing the performance of object detection models. From this review, a diffusion-based text-to-image model, Stable Diffusion XL (SDXL), was selected for synthetic image generation due to its ability to produce visually similar images to the real training set by fine-tuning and to generate diverse images using textual prompts.

For implementation, a pipeline was developed for synthetic image generation and object detection model training. In the first stage, a LoRA model was trained to fine-tune SDXL using 31 images captured with a near-infrared sensor. For LoRA training, the images were resized to match the SDXL model's requirements, and descriptive captions were created for each image. By combining the trained LoRA with the SDXL model, it was possible to generate synthetic images resembling the near-infrared images. The generation of diverse images was automated using the Wildcard extension tool, which allowed controlled randomization of prompts. After filtering and cleaning the outputs, a total of 2500 synthetic images were obtained and used for training the object detection model.

To evaluate the impact of synthetic images, a total of 23 YOLOv8n object detection models were trained with identical training parameters using different combinations of real and synthetic images. Quantitatively, models fine-tuned solely on synthetic images achieved AP\textsubscript{IoU:0.50} between 0.866 and 0.908, outperforming the base model at 0.815. Adding real images on top of a large synthetic set further improved both accuracy and stability. The best performing model combined 568 real images with 2500 synthetic images and achieved an AP\textsubscript{IoU:0.50} of 0.938.

In the qualitative evaluation, where the detections of the base model were compared with best-performing fine-tuned model trained on both synthetic and real images, it was observed that the mixed-dataset model had become substantially better at detecting people in near-infrared images. This improvement was seen not only in simple cases but also in more challenging ones, including varying lighting, partial occlusions and unusual poses of individuals.

The methods implemented in this work can be considered an effective way to expand the training dataset for object detection models while also improving model performance. The approach reduces the need to collect large amounts of real data and enables the creation of more diverse training datasets. To make the method even more efficient and time-saving, future research should focus on automating the annotation process of synthetic images, enabling the generated images to be used directly for model training.

